\section{Annual BAU Comparisons and Reference Diagnostics}
\label{app:annual-bau}

The paired \gls{bau} comparison in Section~\ref{sec:control-common-protocol} covers Core-30 Normal, Topology, Safety, Robustness, and Combined, together with Premium-100 Normal and All-in. Each scenario has a reference without stationary storage and a reference using storage for local \gls{pv} self-consumption. The 2023 and 2024 accounts remain separate; the 2024 reference accompanies the principal controller comparison.

\subsection{Development-Year Annual Comparison}

\noindent\textcolor{orange!70!black}{\textbf{Illustrative annual comparison: invented values.}} The two tables below describe the intended full-year comparison for 2023. Their numbers are manuscript examples, independent of the real Core-30 diagnostic reported later in this section. The chosen year-to-year differences carry no empirical interpretation.

\begingroup\footnotesize
\setlength{\tabcolsep}{3pt}
\renewcommand{\arraystretch}{1.08}
\begin{longtable}{@{}lrrrrrrrr@{}}
\caption{\illustrativecaption Energy and economic outcomes for the two \gls{bau} variants over the complete 2023 window.}
\label{tab:control-bau-2023-energy}\\
\toprule
\textbf{\gls{bau} variant} & \textbf{Cost} & \textbf{CO$_2$} & \textbf{Import} & \textbf{Export} & \textbf{Solar} & \textbf{Daily} & \textbf{Max.} & \textbf{Ramp}\\
 & [EUR] & [kg] & [kWh] & [kWh] & [\%] & [kW] & [kW] & [kW]\\
\midrule
\endfirsthead
\caption[]{\illustrativecaption Energy and economic outcomes for the two \gls{bau} variants over the complete 2023 window. (continued).}\\
\toprule
\textbf{\gls{bau} variant} & \textbf{Cost} & \textbf{CO$_2$} & \textbf{Import} & \textbf{Export} & \textbf{Solar} & \textbf{Daily} & \textbf{Max.} & \textbf{Ramp}\\
 & [EUR] & [kg] & [kWh] & [kWh] & [\%] & [kW] & [kW] & [kW]\\
\midrule
\endhead
\midrule
\multicolumn{9}{r}{\emph{Continued on the next page.}}\\
\endfoot
\bottomrule
\endlastfoot
\multicolumn{9}{@{}l}{\emph{Core-30 Normal}}\\*
Without storage & 29\,160 & 37\,948 & 205\,965 & 42\,840 & 30.3 & 65.9 & 112.7 & 200\,720\\*
With storage & 25\,920 & 33\,920 & 186\,300 & 21\,420 & 65.2 & 61.8 & 107.6 & 187\,200\\
\addlinespace[3pt]
\multicolumn{9}{@{}l}{\emph{Core-30 Topology}}\\*
Without storage & 27\,432 & 36\,040 & 196\,650 & 39\,780 & 30.7 & 62.8 & 110.7 & 192\,400\\*
With storage & 24\,624 & 32\,542 & 178\,848 & 23\,970 & 58.2 & 58.7 & 104.5 & 180\,960\\
\addlinespace[3pt]
\multicolumn{9}{@{}l}{\emph{Core-30 Safety}}\\*
Without storage & 28\,728 & 37\,418 & 203\,895 & 42\,840 & 30.3 & 63.9 & 104.5 & 191\,360\\*
With storage & 25\,596 & 33\,496 & 185\,265 & 21\,930 & 64.3 & 59.7 & 100.4 & 179\,920\\
\addlinespace[3pt]
\multicolumn{9}{@{}l}{\emph{Core-30 Robustness}}\\*
Without storage & 29\,376 & 38\,160 & 207\,000 & 43\,350 & 29.5 & 67.0 & 115.8 & 212\,160\\*
With storage & 26\,568 & 34\,768 & 189\,922 & 24\,990 & 59.4 & 63.9 & 112.7 & 201\,760\\
\addlinespace[3pt]
\multicolumn{9}{@{}l}{\emph{Core-30 Combined}}\\*
Without storage & 27\,540 & 35\,934 & 195\,615 & 39\,984 & 30.3 & 61.8 & 102.5 & 193\,440\\*
With storage & 25\,056 & 33\,072 & 180\,090 & 26\,316 & 54.2 & 58.7 & 99.4 & 185\,120\\
\addlinespace[3pt]
\multicolumn{9}{@{}l}{\emph{Premium-100 Normal}}\\*
Without storage & 119\,880 & 154\,760 & 807\,300 & 295\,800 & 38.6 & 262.7 & 502.2 & 780\,000\\*
With storage & 103\,680 & 135\,680 & 705\,870 & 214\,200 & 55.5 & 236.9 & 456.1 & 717\,600\\
\addlinespace[3pt]
\multicolumn{9}{@{}l}{\emph{Premium-100 All-in}}\\*
Without storage & 113\,400 & 148\,400 & 755\,550 & 269\,280 & 37.4 & 253.4 & 471.5 & 796\,640\\*
With storage & 100\,440 & 132\,500 & 673\,785 & 201\,960 & 53.1 & 231.8 & 440.7 & 748\,800\\
\end{longtable}
\noindent Import/export use the external \gls{rec} boundary; cost sums member purchase bills. All-in cost precedes the separate illustrative \gls{dr} settlement in Table~\ref{tab:control-dr-results}. Daily/Max. denote mean daily/maximum peak; Ramp is accumulated upward power change. Each pair uses the same scenario and year; different populations are not ranked by total cost.\par
\endgroup

\begingroup\footnotesize
\setlength{\tabcolsep}{3pt}
\renewcommand{\arraystretch}{1.08}
\begin{longtable}{@{}lrrrrrr@{}}
\caption{\illustrativecaption Departure service, safety, and storage use for the two \gls{bau} variants over the complete 2023 window.}
\label{tab:control-bau-2023-service}\\
\toprule
\textbf{\gls{bau} variant} & \textbf{\gls{ev} min.} & \textbf{Missed \gls{ev}} & \textbf{\gls{ev} band} & \textbf{Viol.} & \textbf{Corrected} & \textbf{Battery}\\
 & [\%] & [$n/N$] & [\%] & [kWh] & [\%] & [kWh]\\
\midrule
\endfirsthead
\caption[]{\illustrativecaption Departure service, safety, and storage use for the two \gls{bau} variants over the complete 2023 window. (continued).}\\
\toprule
\textbf{\gls{bau} variant} & \textbf{\gls{ev} min.} & \textbf{Missed \gls{ev}} & \textbf{\gls{ev} band} & \textbf{Viol.} & \textbf{Corrected} & \textbf{Battery}\\
 & [\%] & [$n/N$] & [\%] & [kWh] & [\%] & [kWh]\\
\midrule
\endhead
\midrule
\multicolumn{7}{r}{\emph{Continued on the next page.}}\\
\endfoot
\bottomrule
\endlastfoot
\multicolumn{7}{@{}l}{\emph{Core-30 Normal}}\\*
Without storage & 100.0 & 0/10\,000 & 18.0 & 0 & 3.4 & 0\\*
With storage & 100.0 & 0/10\,000 & 18.0 & 0 & 3.0 & 22\,440\\
\addlinespace[3pt]
\multicolumn{7}{@{}l}{\emph{Core-30 Topology}}\\*
Without storage & 99.0 & 90/9\,000 & 18.0 & 0 & 4.0 & 0\\*
With storage & 99.2 & 72/9\,000 & 18.2 & 0 & 3.5 & 19\,380\\
\addlinespace[3pt]
\multicolumn{7}{@{}l}{\emph{Core-30 Safety}}\\*
Without storage & 96.0 & 400/10\,000 & 17.0 & 0 & 14.0 & 0\\*
With storage & 97.0 & 300/10\,000 & 17.5 & 0 & 11.0 & 21\,420\\
\addlinespace[3pt]
\multicolumn{7}{@{}l}{\emph{Core-30 Robustness}}\\*
Without storage & 96.5 & 350/10\,000 & 17.0 & 0 & 10.0 & 0\\*
With storage & 97.2 & 280/10\,000 & 17.5 & 0 & 9.0 & 19\,890\\
\addlinespace[3pt]
\multicolumn{7}{@{}l}{\emph{Core-30 Combined}}\\*
Without storage & 93.0 & 630/9\,000 & 16.0 & 0 & 19.0 & 0\\*
With storage & 94.5 & 495/9\,000 & 16.5 & 0 & 16.0 & 17\,850\\
\addlinespace[3pt]
\multicolumn{7}{@{}l}{\emph{Premium-100 Normal}}\\*
Without storage & 99.0 & 300/30\,000 & 18.0 & 0 & 6.0 & 0\\*
With storage & 99.3 & 210/30\,000 & 18.5 & 0 & 4.5 & 193\,800\\
\addlinespace[3pt]
\multicolumn{7}{@{}l}{\emph{Premium-100 All-in}}\\*
Without storage & 91.0 & 2\,250/25\,000 & 16.0 & 0 & 23.0 & 0\\*
With storage & 93.0 & 1\,750/25\,000 & 16.5 & 0 & 19.0 & 153\,000\\
\end{longtable}
\noindent Missed \gls{ev} gives failed minimum-service departures over an invented common eligible cohort within each pair. These counts are invented cohorts, independent of the dataset audit. \gls{v2g} is disabled. Battery denotes stationary-battery throughput. Appliance, \gls{dr}, and runtime event outcomes are illustrated separately in Tables~\ref{tab:control-runtime-events} and~\ref{tab:control-dr-results}.\par
\endgroup


The illustrative tables use common invented eligible-departure counts within each scenario pair: 10\,000 in Core-30 Normal, Safety, and Robustness; 9\,000 in Topology and Combined; 30\,000 in Premium Normal; and 25\,000 in All-in. They provide arithmetically consistent examples of percentages and missed-service counts. These are not counts from the dataset audit. The annual controller fixtures add an invented unattainable cohort equal to 2\% of eligible departures and assume zero incomplete records. These additional requests remain outside the service-rate denominator. Actual comparisons retain both scenario-specific eligibility and the common cohort required for paired disturbance assessment.

\subsection{Companion Service and Execution Records}

Service, runtime, computation, and deployment accompany the annual scorecards in Tables~\ref{tab:control-safety-events}, \ref{tab:control-robustness-events}, \ref{tab:control-runtime-events}, \ref{tab:control-computation}, and~\ref{tab:control-deployment}. The Premium All-in response account appears in Table~\ref{tab:control-dr-results}. Each companion uses explicitly invented counts and magnitudes. \gls{v2g} is disabled, stationary-battery throughput is zero without storage, and \gls{bau} with storage has zero members saving against itself.

\subsection{Real Short-Window Reference Diagnostics}
\label{app:reference-diagnostics}

The following values were obtained from an actual Core-30 simulator execution. They preserve its original window and accounting. The annual illustrative tables above are independent manuscript examples.

\begin{table}[H]
\centering
\caption{\textbf{Measured short-window \gls{bau} diagnostic.} Portugal v3.1, 671 executed fifteen-minute steps from 1 January 2023 00:00 through 7 January 2023 23:30. Cost and import sum member-level purchases; peak uses \gls{rec} net power.}
\label{tab:control-real-bau}
\small
\begin{tabular}{@{}lrrrrrr@{}}
\toprule
\textbf{Dataset} & \textbf{Cost} & \textbf{Import} & \textbf{Peak} & \textbf{\gls{ev}} & \textbf{Min.} & \textbf{Band}\\
 & [EUR] & [kWh] & [kW] & departures & [\%] & [\%]\\
\midrule
Core-30 & 240.59 & 5074.2 & 85.59 & 105 & 100.0 & 6.7 \\
\bottomrule
\end{tabular}
\end{table}

The storage-enabled \gls{bau} execution in Table~\ref{tab:control-real-bau} provides a concrete reference on the current Core-30 inputs. All 105 observed departures were eligible for minimum-service evaluation, and every one met that criterion. The target-band rate was much lower: 7 of 105 departures. Immediate charging towards full charge explains why ordinary operation can satisfy the minimum requirement while delivering more charge than requested. The run recorded zero residual electrical-violation energy under the simulator's constraint handling. Its exported window contains 671 action intervals; no unexecuted final interval is added to the reported energy or cost.

The import costs were calculated from each member's exported positive net energy and the corresponding realised tariff, retaining negative prices and zero export remuneration. This run establishes actual short-window cost and service values; the full test-window comparison retains its separate scope. It also illustrates why a low bill, a safe executed trajectory, and accurate service are separate observations. The annual scorecard retains these separate outcomes.

The storage-enabled diagnostic used a state-of-charge range of 0.2--0.9, charge and discharge command rates of 0.35, and a local import threshold of 2~kW for discharge. \gls{ev} charging used rate 1.0 and target \gls{soc} 1.0; the appliance start command was 1.0. The no-storage comparison sets stationary-battery charge and discharge to zero while retaining the remaining rules. The resolved configuration is archived with each execution.

\subsection{Complete Annual Controller Scorecards}
\label{app:control-variant-coverage}

The principal controller families and the additional training and hierarchy variants share the seven-scenario scope. Table~\ref{tab:control-variant-coverage} first compares the annual cost of the additional configurations. The subsequent paired tables give energy, service, safety, resource use, and member benefit for all 26 configurations. Every value is an invented editorial fixture. The complete-year account includes supervisory fallback when the hypothetical learned controller cannot continue; it does not assert that each implementation has been executed on every scenario.

\begin{table}[!htbp]
\centering
\caption{\illustrativecaption Annual energy-purchase cost [EUR] of the additional training and hierarchy variants.}
\label{tab:control-variant-coverage}
\footnotesize
\setlength{\tabcolsep}{2.5pt}
\renewcommand{\arraystretch}{1.10}
\begin{tabularx}{\textwidth}{@{}p{0.28\textwidth}*{7}{>{\raggedleft\arraybackslash}X}@{}}
\toprule
Configuration & C30 N & C30 T & C30 S & C30 R & C30 C & P100 N & P100 A\\
\midrule
\gls{matd3} unguided & 21\,840 & 22\,408 & 21\,998 & 25\,520 & 25\,123 & 87\,360 & 104\,095\\
\gls{matd3} guided exploration & 20\,880 & 21\,423 & 21\,031 & 24\,398 & 24\,019 & 83\,520 & 99\,519\\
\gls{matd3} demonstrations & 20\,400 & 20\,930 & 20\,548 & 23\,837 & 23\,467 & 81\,600 & 97\,232\\
\gls{rbc} Smart Local global & 20\,880 & 20\,332 & 21\,031 & 22\,686 & 22\,202 & 83\,520 & 91\,428\\
\gls{rbc} Smart Local per-member & 20\,640 & 20\,098 & 20\,790 & 22\,425 & 21\,947 & 82\,560 & 90\,377\\
\gls{matd3} residual global & 19\,920 & 20\,438 & 20\,064 & 23\,277 & 22\,915 & 79\,680 & 94\,944\\
\gls{matd3} residual per-member & 20\,160 & 20\,684 & 20\,306 & 23\,557 & 23\,191 & 80\,640 & 96\,088\\
\bottomrule
\end{tabularx}
\par\smallskip\begin{minipage}{\textwidth}\footnotesize C30: Core-30; P100: Premium-100. N: Normal; T: Topology; S: Safety; R: Robustness; C: Combined; A: All-in. Each column covers its complete 2024 sequence, including supervisory fallback. *Generic offline-policy placeholder; the associated method remains to be specified.\end{minipage}
\end{table}

\noindent\textbf{All values in the following scorecards are invented manuscript fixtures.} Each scenario covers a complete 2024 comparison. Codes retain the same controller configuration throughout the appendix. Operational and computational companions use the same explicitly hypothetical status.
\par Solar denotes \gls{pv} self-consumption; Daily and Max. denote \gls{rec} peaks. Saving counts members with a lower bill than B1. The invented eligible-departure cohorts follow the \gls{bau} tables; unattainable requests remain separate. \gls{v2g} is disabled.
\par\medskip\noindent\begin{minipage}{\textwidth}\singlespacing\footnotesize\begin{tabularx}{\textwidth}{@{}lXlX@{}}\toprule Code & Configuration & Code & Configuration\\\midrule
B0 & \gls{bau} without storage & MP & \gls{mappo}\\
B1 & \gls{bau} with storage & HP & \gls{happo}\\
RB & \gls{rbc} Basic & MS & \gls{masac}\\
RSL & \gls{rbc} Smart Local & OF & Offline policy*\\
RC & \gls{rbc} Community & TP & Transformer \gls{ppo}\\
MLF & \gls{milp} local fixed & TT & Transformer \gls{matd3}\\
MCF & \gls{milp} \gls{rec} fixed & TU & \gls{matd3} unguided\\
MLA & \gls{milp} local flexible & TG & \gls{matd3} guided exploration\\
MCA & \gls{milp} \gls{rec} flexible & TM & \gls{matd3} demonstrations\\
LP & Local \gls{ppo} & RG & \gls{rbc} Smart Local global\\
LT & Local \gls{td3} & RM & \gls{rbc} Smart Local per-member\\
DD & \gls{maddpg} & HG & \gls{matd3} residual global\\
TD & \gls{matd3} (residual) & HM & \gls{matd3} residual per-member\\
\bottomrule\end{tabularx}\end{minipage}
\clearpage
\noindent\begin{minipage}{\textwidth}\singlespacing\captionsetup{font=footnotesize}
\setlength{\tabcolsep}{4pt}
\renewcommand{\arraystretch}{1.0}
\captionof{table}{\illustrativecaption Core-30 Normal. Energy and economic outcomes.}
\label{tab:full-scorecard-0-0}
\fontsize{8}{9}\selectfont
\begin{tabular*}{\textwidth}{@{\extracolsep{\fill}}lrrrrrrrr@{}}
\toprule
Code & Cost & CO$_2$ & Import & Export & Solar & Daily & Max. & Ramp\\
 & [EUR] & [kg] & [kWh] & [kWh] & [\%] & [kW] & [kW] & [kW]\\
\midrule
B0 & 27\,000 & 35\,800 & 199\,000 & 42\,000 & 30.0 & 64.0 & 110.0 & 193\,000\\
B1 & 24\,000 & 32\,000 & 180\,000 & 21\,000 & 65.0 & 60.0 & 105.0 & 180\,000\\
RB & 23\,040 & 30\,800 & 175\,000 & 20\,400 & 66.0 & 57.0 & 100.0 & 174\,600\\
RSL & 21\,600 & 28\,400 & 166\,000 & 17\,400 & 71.0 & 54.0 & 96.0 & 167\,400\\
RC & 21\,120 & 27\,700 & 163\,000 & 16\,800 & 72.0 & 51.6 & 91.0 & 158\,400\\
MLF & 21\,480 & 27\,505 & 169\,177 & 15\,717 & 73.8 & 53.7 & 93.9 & 166\,994\\
MCF & 21\,000 & 26\,890 & 165\,396 & 15\,366 & 74.4 & 52.5 & 91.8 & 163\,262\\
MLA & 19\,560 & 25\,046 & 154\,055 & 14\,312 & 76.1 & 48.9 & 85.5 & 152\,067\\
MCA & 18\,720 & 23\,971 & 147\,439 & 13\,698 & 77.2 & 46.8 & 81.8 & 145\,537\\
LP & 22\,080 & 29\,200 & 169\,000 & 19\,200 & 68.0 & 56.4 & 99.0 & 172\,800\\
LT & 21\,600 & 28\,800 & 167\,000 & 18\,600 & 69.0 & 55.2 & 97.0 & 169\,200\\
DD & 20\,880 & 27\,800 & 163\,500 & 17\,400 & 71.0 & 52.8 & 94.0 & 162\,000\\
TD & 19\,680 & 25\,200 & 155\,000 & 14\,400 & 76.0 & 49.2 & 86.0 & 153\,000\\
MP & 20\,640 & 26\,900 & 161\,000 & 16\,200 & 73.0 & 51.6 & 92.0 & 160\,200\\
HP & 20\,400 & 26\,500 & 159\,500 & 15\,600 & 74.0 & 51.0 & 90.0 & 158\,400\\
MS & 20\,640 & 27\,100 & 162\,000 & 16\,800 & 72.0 & 52.2 & 93.0 & 162\,000\\
OF & 21\,240 & 28\,800 & 167\,000 & 18\,600 & 69.0 & 55.2 & 97.0 & 169\,200\\
TP & 21\,360 & 28\,100 & 164\,500 & 18\,000 & 70.0 & 54.0 & 95.0 & 167\,400\\
TT & 20\,160 & 25\,700 & 157\,000 & 15\,000 & 75.0 & 50.4 & 88.0 & 156\,600\\
TU & 21\,840 & 27\,966 & 172\,012 & 15\,980 & 73.4 & 54.6 & 95.4 & 169\,793\\
TG & 20\,880 & 26\,737 & 164\,451 & 15\,278 & 74.5 & 52.2 & 91.2 & 162\,329\\
TM & 20\,400 & 26\,122 & 160\,671 & 14\,927 & 75.1 & 51.0 & 89.1 & 158\,598\\
RG & 20\,880 & 27\,453 & 160\,467 & 16\,200 & 73.0 & 51.0 & 92.8 & 158\,400\\
RM & 20\,640 & 27\,138 & 158\,622 & 15\,600 & 74.0 & 49.8 & 91.7 & 156\,600\\
HG & 19\,920 & 25\,507 & 156\,890 & 13\,800 & 77.0 & 46.8 & 87.0 & 144\,000\\
HM & 20\,160 & 25\,815 & 158\,780 & 13\,200 & 78.0 & 45.6 & 88.1 & 140\,400\\
\bottomrule\end{tabular*}
\end{minipage}\par\medskip
\noindent\begin{minipage}{\textwidth}\singlespacing\captionsetup{font=footnotesize}
\setlength{\tabcolsep}{4pt}
\renewcommand{\arraystretch}{1.0}
\captionof{table}{\illustrativecaption Core-30 Normal. Service, safety and resource use.}
\label{tab:full-scorecard-0-1}
\fontsize{8}{9}\selectfont
\begin{tabular*}{\textwidth}{@{\extracolsep{\fill}}lrrrrrrr@{}}
\toprule
Code & \gls{ev} min. & \gls{ev} band & Viol. & Corr. & Battery & \gls{v2g} & Saving\\
 & [\%] & [\%] & [kWh] & [\%] & [kWh] & [kWh] & [n/N]\\
\midrule
B0 & 100.0 & 18.0 & 0 & 3.4 & 0 & 0 & 0/30\\
B1 & 100.0 & 18.0 & 0 & 3.0 & 22\,000 & 0 & 0/30\\
RB & 99.2 & 75.0 & 0 & 7.0 & 23\,500 & 0 & 21/30\\
RSL & 99.8 & 92.0 & 0 & 4.0 & 25\,500 & 0 & 26/30\\
RC & 99.8 & 93.0 & 0 & 3.6 & 26\,800 & 0 & 27/30\\
MLF & 100.0 & 18.0 & 0 & 1.2 & 30\,561 & 0 & 29/30\\
MCF & 100.0 & 18.0 & 0 & 1.2 & 29\,878 & 0 & 29/30\\
MLA & 100.0 & 98.0 & 0 & 1.2 & 27\,829 & 0 & 29/30\\
MCA & 100.0 & 98.0 & 0 & 1.2 & 26\,634 & 0 & 29/30\\
LP & 98.2 & 81.0 & 0 & 9.5 & 28\,200 & 0 & 23/30\\
LT & 98.8 & 85.0 & 0 & 8.0 & 29\,100 & 0 & 25/30\\
DD & 98.6 & 84.0 & 0 & 14.0 & 32\,000 & 0 & 26/30\\
TD & 99.8 & 96.0 & 0 & 2.3 & 28\,000 & 0 & 29/30\\
MP & 99.5 & 91.0 & 0 & 5.5 & 29\,700 & 0 & 27/30\\
HP & 99.6 & 92.5 & 0 & 4.8 & 30\,300 & 0 & 28/30\\
MS & 99.3 & 90.0 & 0 & 6.5 & 31\,000 & 0 & 27/30\\
OF & 99.2 & 89.0 & 0 & 5.8 & 29\,100 & 0 & 25/30\\
TP & 99.0 & 89.0 & 0 & 5.0 & 28\,500 & 0 & 26/30\\
TT & 99.7 & 95.5 & 0 & 2.0 & 27\,600 & 0 & 28/30\\
TU & 96.0 & 73.0 & 0 & 12.4 & 33\,700 & 0 & 29/30\\
TG & 98.1 & 85.0 & 0 & 8.2 & 31\,300 & 0 & 29/30\\
TM & 99.0 & 90.0 & 0 & 6.1 & 29\,400 & 0 & 29/30\\
RG & 99.8 & 93.0 & 0 & 3.8 & 24\,650 & 0 & 28/30\\
RM & 99.9 & 94.0 & 0 & 3.5 & 24\,367 & 0 & 29/30\\
HG & 99.8 & 96.0 & 0 & 2.5 & 28\,341 & 0 & 28/30\\
HM & 99.8 & 96.0 & 0 & 2.7 & 28\,683 & 0 & 27/30\\
\bottomrule\end{tabular*}
\end{minipage}\par\medskip
\clearpage
\noindent\begin{minipage}{\textwidth}\singlespacing\captionsetup{font=footnotesize}
\setlength{\tabcolsep}{4pt}
\renewcommand{\arraystretch}{1.0}
\captionof{table}{\illustrativecaption Core-30 Topology. Energy and economic outcomes.}
\label{tab:full-scorecard-1-0}
\fontsize{8}{9}\selectfont
\begin{tabular*}{\textwidth}{@{\extracolsep{\fill}}lrrrrrrrr@{}}
\toprule
Code & Cost & CO$_2$ & Import & Export & Solar & Daily & Max. & Ramp\\
 & [EUR] & [kg] & [kWh] & [kWh] & [\%] & [kW] & [kW] & [kW]\\
\midrule
B0 & 25\,400 & 34\,000 & 190\,000 & 39\,000 & 30.4 & 61.0 & 108.0 & 185\,000\\
B1 & 22\,800 & 30\,700 & 172\,800 & 23\,500 & 58.0 & 57.0 & 102.0 & 174\,000\\
RB & 23\,639 & 29\,549 & 168\,000 & 22\,829 & 59.2 & 54.1 & 97.1 & 168\,780\\
RSL & 21\,033 & 27\,246 & 159\,360 & 19\,471 & 65.2 & 51.3 & 93.3 & 161\,820\\
RC & 20\,566 & 26\,575 & 156\,480 & 18\,800 & 66.4 & 49.0 & 88.4 & 153\,120\\
MLF & 20\,610 & 26\,388 & 162\,410 & 17\,588 & 68.6 & 51.0 & 91.2 & 161\,428\\
MCF & 20\,150 & 25\,798 & 158\,780 & 17\,195 & 69.3 & 49.9 & 89.2 & 157\,820\\
MLA & 18\,768 & 24\,029 & 147\,893 & 16\,016 & 71.4 & 46.5 & 83.1 & 146\,998\\
MCA & 17\,962 & 22\,997 & 141\,541 & 15\,329 & 72.6 & 44.5 & 79.5 & 140\,686\\
LP & 22\,654 & 28\,014 & 162\,240 & 21\,486 & 61.6 & 53.6 & 96.2 & 167\,040\\
LT & 22\,162 & 27\,630 & 160\,320 & 20\,814 & 62.8 & 52.4 & 94.2 & 163\,560\\
DD & 21\,423 & 26\,671 & 156\,960 & 19\,471 & 65.2 & 50.2 & 91.3 & 156\,600\\
TD & 20\,192 & 24\,176 & 148\,800 & 16\,114 & 71.2 & 46.7 & 83.5 & 147\,900\\
MP & 21\,177 & 25\,807 & 154\,560 & 18\,129 & 67.6 & 49.0 & 89.4 & 154\,860\\
HP & 20\,930 & 25\,423 & 153\,120 & 17\,457 & 68.8 & 48.5 & 87.4 & 153\,120\\
MS & 21\,177 & 25\,999 & 155\,520 & 18\,800 & 66.4 & 49.6 & 90.3 & 156\,600\\
OF & 21\,792 & 27\,630 & 160\,320 & 20\,814 & 62.8 & 52.4 & 94.2 & 163\,560\\
TP & 20\,799 & 26\,958 & 157\,920 & 20\,143 & 64.0 & 51.3 & 92.3 & 161\,820\\
TT & 19\,631 & 24\,656 & 150\,720 & 16\,786 & 70.0 & 47.9 & 85.5 & 151\,380\\
TU & 22\,408 & 26\,830 & 165\,132 & 17\,882 & 68.1 & 51.9 & 92.7 & 164\,133\\
TG & 21\,423 & 25\,651 & 157\,873 & 17\,097 & 69.5 & 49.6 & 88.6 & 156\,918\\
TM & 20\,930 & 25\,061 & 154\,244 & 16\,704 & 70.2 & 48.5 & 86.6 & 153\,311\\
RG & 20\,332 & 26\,338 & 154\,048 & 18\,129 & 67.6 & 48.5 & 90.1 & 153\,120\\
RM & 20\,098 & 26\,036 & 152\,277 & 17\,457 & 68.8 & 47.3 & 89.1 & 151\,380\\
HG & 20\,438 & 24\,471 & 150\,614 & 15\,443 & 72.4 & 44.5 & 84.5 & 139\,200\\
HM & 20\,684 & 24\,766 & 152\,429 & 14\,771 & 73.6 & 43.3 & 85.6 & 135\,720\\
\bottomrule\end{tabular*}
\end{minipage}\par\medskip
\noindent\begin{minipage}{\textwidth}\singlespacing\captionsetup{font=footnotesize}
\setlength{\tabcolsep}{4pt}
\renewcommand{\arraystretch}{1.0}
\captionof{table}{\illustrativecaption Core-30 Topology. Service, safety and resource use.}
\label{tab:full-scorecard-1-1}
\fontsize{8}{9}\selectfont
\begin{tabular*}{\textwidth}{@{\extracolsep{\fill}}lrrrrrrr@{}}
\toprule
Code & \gls{ev} min. & \gls{ev} band & Viol. & Corr. & Battery & \gls{v2g} & Saving\\
 & [\%] & [\%] & [kWh] & [\%] & [kWh] & [kWh] & [n/N]\\
\midrule
B0 & 99.0 & 18.0 & 0 & 4.0 & 0 & 0 & 0/30\\
B1 & 99.2 & 18.2 & 0 & 3.5 & 19\,000 & 0 & 0/30\\
RB & 96.2 & 69.6 & 0 & 10.0 & 20\,295 & 0 & 7/30\\
RSL & 99.0 & 90.6 & 0 & 7.0 & 22\,023 & 0 & 28/30\\
RC & 99.0 & 91.6 & 0 & 6.6 & 23\,145 & 0 & 28/30\\
MLF & 99.8 & 17.6 & 0 & 4.2 & 26\,394 & 0 & 28/30\\
MCF & 99.8 & 17.6 & 0 & 4.2 & 25\,804 & 0 & 28/30\\
MLA & 99.8 & 97.6 & 0 & 4.2 & 24\,034 & 0 & 28/30\\
MCA & 99.8 & 97.6 & 0 & 4.2 & 23\,002 & 0 & 28/30\\
LP & 95.2 & 75.6 & 0 & 12.5 & 24\,355 & 0 & 28/30\\
LT & 95.8 & 79.6 & 0 & 11.0 & 25\,132 & 0 & 28/30\\
DD & 95.6 & 78.6 & 0 & 17.0 & 27\,636 & 0 & 28/30\\
TD & 96.8 & 90.6 & 0 & 5.3 & 24\,182 & 0 & 28/30\\
MP & 96.5 & 85.6 & 0 & 8.5 & 25\,650 & 0 & 28/30\\
HP & 96.6 & 87.1 & 0 & 7.8 & 26\,168 & 0 & 28/30\\
MS & 96.3 & 84.6 & 0 & 9.5 & 26\,773 & 0 & 28/30\\
OF & 96.2 & 83.6 & 0 & 8.8 & 25\,132 & 0 & 28/30\\
TP & 98.2 & 87.6 & 0 & 8.0 & 24\,614 & 0 & 28/30\\
TT & 98.9 & 94.1 & 0 & 5.0 & 23\,836 & 0 & 28/30\\
TU & 93.0 & 67.6 & 0 & 15.4 & 29\,105 & 0 & 28/30\\
TG & 95.1 & 79.6 & 0 & 11.2 & 27\,032 & 0 & 28/30\\
TM & 96.0 & 84.6 & 0 & 9.1 & 25\,391 & 0 & 28/30\\
RG & 99.0 & 91.6 & 0 & 6.8 & 21\,289 & 0 & 28/30\\
RM & 99.1 & 92.6 & 0 & 6.5 & 21\,044 & 0 & 28/30\\
HG & 96.8 & 90.6 & 0 & 5.5 & 24\,476 & 0 & 28/30\\
HM & 96.8 & 90.6 & 0 & 5.7 & 24\,772 & 0 & 28/30\\
\bottomrule\end{tabular*}
\end{minipage}\par\medskip
\clearpage
\noindent\begin{minipage}{\textwidth}\singlespacing\captionsetup{font=footnotesize}
\setlength{\tabcolsep}{4pt}
\renewcommand{\arraystretch}{1.0}
\captionof{table}{\illustrativecaption Core-30 Safety. Energy and economic outcomes.}
\label{tab:full-scorecard-2-0}
\fontsize{8}{9}\selectfont
\begin{tabular*}{\textwidth}{@{\extracolsep{\fill}}lrrrrrrrr@{}}
\toprule
Code & Cost & CO$_2$ & Import & Export & Solar & Daily & Max. & Ramp\\
 & [EUR] & [kg] & [kWh] & [kWh] & [\%] & [kW] & [kW] & [kW]\\
\midrule
B0 & 26\,600 & 35\,300 & 197\,000 & 42\,000 & 30.0 & 62.0 & 102.0 & 184\,000\\
B1 & 23\,700 & 31\,600 & 179\,000 & 21\,500 & 64.2 & 58.0 & 98.0 & 173\,000\\
RB & 23\,207 & 30\,415 & 174\,028 & 20\,886 & 65.2 & 55.1 & 93.3 & 167\,810\\
RSL & 21\,757 & 28\,045 & 165\,078 & 17\,814 & 70.3 & 52.2 & 89.6 & 160\,890\\
RC & 21\,273 & 27\,354 & 162\,094 & 17\,200 & 71.3 & 49.9 & 84.9 & 152\,240\\
MLF & 21\,318 & 27\,161 & 168\,237 & 16\,091 & 73.2 & 51.9 & 87.6 & 160\,500\\
MCF & 20\,841 & 26\,554 & 164\,477 & 15\,732 & 73.8 & 50.8 & 85.7 & 156\,913\\
MLA & 19\,412 & 24\,733 & 153\,199 & 14\,653 & 75.6 & 47.3 & 79.8 & 146\,153\\
MCA & 18\,578 & 23\,671 & 146\,620 & 14\,024 & 76.6 & 45.2 & 76.3 & 139\,877\\
LP & 22\,240 & 28\,835 & 168\,061 & 19\,657 & 67.2 & 54.5 & 92.4 & 166\,080\\
LT & 21\,757 & 28\,440 & 166\,072 & 19\,043 & 68.3 & 53.4 & 90.5 & 162\,620\\
DD & 21\,031 & 27\,452 & 162\,592 & 17\,814 & 70.3 & 51.0 & 87.7 & 155\,700\\
TD & 19\,823 & 24\,885 & 154\,139 & 14\,743 & 75.4 & 47.6 & 80.3 & 147\,050\\
MP & 20\,790 & 26\,564 & 160\,106 & 16\,586 & 72.4 & 49.9 & 85.9 & 153\,970\\
HP & 20\,548 & 26\,169 & 158\,614 & 15\,971 & 73.4 & 49.3 & 84.0 & 152\,240\\
MS & 20\,790 & 26\,761 & 161\,100 & 17\,200 & 71.3 & 50.5 & 86.8 & 155\,700\\
OF & 21\,394 & 28\,440 & 166\,072 & 19\,043 & 68.3 & 53.4 & 90.5 & 162\,620\\
TP & 21\,515 & 27\,749 & 163\,586 & 18\,429 & 69.3 & 52.2 & 88.7 & 160\,890\\
TT & 20\,306 & 25\,379 & 156\,128 & 15\,357 & 74.4 & 48.7 & 82.1 & 150\,510\\
TU & 21\,998 & 27\,616 & 171\,056 & 16\,360 & 72.7 & 52.8 & 89.0 & 163\,190\\
TG & 21\,031 & 26\,403 & 163\,537 & 15\,642 & 73.9 & 50.5 & 85.1 & 156\,016\\
TM & 20\,548 & 25\,795 & 159\,778 & 15\,282 & 74.5 & 49.3 & 83.2 & 152\,430\\
RG & 21\,031 & 27\,110 & 159\,576 & 16\,586 & 72.4 & 49.3 & 86.6 & 152\,240\\
RM & 20\,790 & 26\,799 & 157\,741 & 15\,971 & 73.4 & 48.1 & 85.6 & 150\,510\\
HG & 20\,064 & 25\,188 & 156\,018 & 14\,129 & 76.5 & 45.2 & 81.2 & 138\,400\\
HM & 20\,306 & 25\,492 & 157\,898 & 13\,514 & 77.5 & 44.1 & 82.2 & 134\,940\\
\bottomrule\end{tabular*}
\end{minipage}\par\medskip
\noindent\begin{minipage}{\textwidth}\singlespacing\captionsetup{font=footnotesize}
\setlength{\tabcolsep}{4pt}
\renewcommand{\arraystretch}{1.0}
\captionof{table}{\illustrativecaption Core-30 Safety. Service, safety and resource use.}
\label{tab:full-scorecard-2-1}
\fontsize{8}{9}\selectfont
\begin{tabular*}{\textwidth}{@{\extracolsep{\fill}}lrrrrrrr@{}}
\toprule
Code & \gls{ev} min. & \gls{ev} band & Viol. & Corr. & Battery & \gls{v2g} & Saving\\
 & [\%] & [\%] & [kWh] & [\%] & [kWh] & [kWh] & [n/N]\\
\midrule
B0 & 96.0 & 17.0 & 0 & 14.0 & 0 & 0 & 0/30\\
B1 & 97.0 & 17.5 & 0 & 11.0 & 21\,000 & 0 & 0/30\\
RB & 97.2 & 71.4 & 0 & 13.0 & 22\,432 & 0 & 28/30\\
RSL & 97.4 & 88.0 & 0 & 9.0 & 24\,341 & 0 & 28/30\\
RC & 97.8 & 89.4 & 0 & 9.6 & 25\,582 & 0 & 28/30\\
MLF & 99.5 & 17.1 & 0 & 7.2 & 29\,172 & 0 & 28/30\\
MCF & 99.5 & 17.1 & 0 & 7.2 & 28\,520 & 0 & 28/30\\
MLA & 99.5 & 97.1 & 0 & 7.2 & 26\,564 & 0 & 28/30\\
MCA & 99.5 & 97.1 & 0 & 7.2 & 25\,423 & 0 & 28/30\\
LP & 96.2 & 77.4 & 0 & 15.5 & 26\,918 & 0 & 28/30\\
LT & 96.8 & 81.4 & 0 & 14.0 & 27\,777 & 0 & 28/30\\
DD & 92.0 & 74.0 & 0 & 25.0 & 30\,545 & 0 & 28/30\\
TD & 98.6 & 93.0 & 0 & 6.8 & 26\,727 & 0 & 28/30\\
MP & 97.5 & 87.4 & 0 & 11.5 & 28\,350 & 0 & 28/30\\
HP & 97.6 & 88.9 & 0 & 10.8 & 28\,923 & 0 & 28/30\\
MS & 97.3 & 86.4 & 0 & 12.5 & 29\,591 & 0 & 28/30\\
OF & 97.2 & 85.4 & 0 & 11.8 & 27\,777 & 0 & 28/30\\
TP & 97.0 & 85.4 & 0 & 11.0 & 27\,205 & 0 & 28/30\\
TT & 98.8 & 94.0 & 0 & 5.1 & 26\,345 & 0 & 28/30\\
TU & 94.0 & 69.4 & 0 & 18.4 & 32\,168 & 0 & 28/30\\
TG & 96.1 & 81.4 & 0 & 14.2 & 29\,877 & 0 & 28/30\\
TM & 97.0 & 86.4 & 0 & 12.1 & 28\,064 & 0 & 28/30\\
RG & 97.8 & 89.4 & 0 & 9.8 & 23\,530 & 0 & 28/30\\
RM & 97.9 & 90.4 & 0 & 9.5 & 23\,259 & 0 & 28/30\\
HG & 97.8 & 92.4 & 0 & 8.5 & 27\,053 & 0 & 28/30\\
HM & 97.8 & 92.4 & 0 & 8.7 & 27\,379 & 0 & 28/30\\
\bottomrule\end{tabular*}
\end{minipage}\par\medskip
\clearpage
\noindent\begin{minipage}{\textwidth}\singlespacing\captionsetup{font=footnotesize}
\setlength{\tabcolsep}{4pt}
\renewcommand{\arraystretch}{1.0}
\captionof{table}{\illustrativecaption Core-30 Robustness. Energy and economic outcomes.}
\label{tab:full-scorecard-3-0}
\fontsize{8}{9}\selectfont
\begin{tabular*}{\textwidth}{@{\extracolsep{\fill}}lrrrrrrrr@{}}
\toprule
Code & Cost & CO$_2$ & Import & Export & Solar & Daily & Max. & Ramp\\
 & [EUR] & [kg] & [kWh] & [kWh] & [\%] & [kW] & [kW] & [kW]\\
\midrule
B0 & 27\,200 & 36\,000 & 200\,000 & 42\,500 & 29.2 & 65.0 & 113.0 & 204\,000\\
B1 & 24\,600 & 32\,800 & 183\,500 & 24\,500 & 59.2 & 62.0 & 110.0 & 194\,000\\
RB & 26\,922 & 31\,570 & 178\,403 & 23\,800 & 60.3 & 58.9 & 104.8 & 188\,180\\
RSL & 23\,520 & 29\,110 & 169\,228 & 20\,300 & 66.2 & 55.8 & 100.6 & 180\,420\\
RC & 22\,947 & 28\,392 & 166\,169 & 19\,600 & 67.3 & 53.3 & 95.3 & 170\,720\\
MLF & 24\,219 & 28\,193 & 172\,467 & 18\,336 & 69.4 & 55.5 & 98.4 & 179\,982\\
MCF & 23\,678 & 27\,562 & 168\,612 & 17\,927 & 70.1 & 54.2 & 96.2 & 175\,960\\
MLA & 22\,054 & 25\,672 & 157\,051 & 16\,697 & 72.2 & 50.5 & 89.6 & 163\,894\\
MCA & 21\,107 & 24\,570 & 150\,306 & 15\,981 & 73.4 & 48.4 & 85.7 & 156\,857\\
LP & 25\,800 & 29\,930 & 172\,286 & 22\,400 & 62.7 & 58.3 & 103.7 & 186\,240\\
LT & 25\,240 & 29\,520 & 170\,247 & 21\,700 & 63.8 & 57.0 & 101.6 & 182\,360\\
DD & 24\,720 & 28\,495 & 166\,679 & 20\,300 & 66.2 & 54.6 & 98.5 & 174\,600\\
TD & 22\,080 & 25\,830 & 158\,014 & 16\,800 & 72.0 & 50.8 & 90.1 & 164\,900\\
MP & 23\,520 & 27\,572 & 164\,131 & 18\,900 & 68.5 & 53.3 & 96.4 & 172\,660\\
HP & 23\,837 & 27\,162 & 162\,601 & 18\,200 & 69.7 & 52.7 & 94.3 & 170\,720\\
MS & 24\,118 & 27\,778 & 165\,150 & 19\,600 & 67.3 & 53.9 & 97.4 & 174\,600\\
OF & 24\,819 & 29\,520 & 170\,247 & 21\,700 & 63.8 & 57.0 & 101.6 & 182\,360\\
TP & 23\,208 & 28\,802 & 167\,699 & 21\,000 & 65.0 & 55.8 & 99.5 & 180\,420\\
TT & 21\,840 & 26\,342 & 160\,053 & 17\,500 & 70.8 & 52.1 & 92.2 & 168\,780\\
TU & 25\,520 & 28\,665 & 175\,357 & 18\,643 & 68.9 & 56.4 & 99.9 & 182\,999\\
TG & 24\,398 & 27\,405 & 167\,649 & 17\,824 & 70.3 & 53.9 & 95.5 & 174\,955\\
TM & 23\,837 & 26\,775 & 163\,795 & 17\,415 & 71.0 & 52.7 & 93.3 & 170\,933\\
RG & 22\,686 & 28\,139 & 163\,587 & 18\,900 & 68.5 & 52.7 & 97.2 & 170\,720\\
RM & 22\,425 & 27\,816 & 161\,706 & 18\,200 & 69.7 & 51.5 & 96.1 & 168\,780\\
HG & 23\,277 & 26\,145 & 159\,941 & 16\,100 & 73.2 & 48.4 & 91.1 & 155\,200\\
HM & 23\,557 & 26\,460 & 161\,867 & 15\,400 & 74.3 & 47.1 & 92.3 & 151\,320\\
\bottomrule\end{tabular*}
\end{minipage}\par\medskip
\noindent\begin{minipage}{\textwidth}\singlespacing\captionsetup{font=footnotesize}
\setlength{\tabcolsep}{4pt}
\renewcommand{\arraystretch}{1.0}
\captionof{table}{\illustrativecaption Core-30 Robustness. Service, safety and resource use.}
\label{tab:full-scorecard-3-1}
\fontsize{8}{9}\selectfont
\begin{tabular*}{\textwidth}{@{\extracolsep{\fill}}lrrrrrrr@{}}
\toprule
Code & \gls{ev} min. & \gls{ev} band & Viol. & Corr. & Battery & \gls{v2g} & Saving\\
 & [\%] & [\%] & [kWh] & [\%] & [kWh] & [kWh] & [n/N]\\
\midrule
B0 & 96.5 & 17.0 & 0 & 10.0 & 0 & 0 & 0/30\\
B1 & 97.2 & 17.5 & 0 & 9.0 & 19\,500 & 0 & 0/30\\
RB & 95.2 & 67.8 & 0 & 15.0 & 20\,830 & 0 & 7/30\\
RSL & 97.7 & 87.0 & 0 & 12.0 & 22\,602 & 0 & 27/30\\
RC & 97.5 & 88.9 & 0 & 11.6 & 23\,755 & 0 & 27/30\\
MLF & 98.6 & 15.5 & 0 & 9.2 & 27\,088 & 0 & 27/30\\
MCF & 98.6 & 15.5 & 0 & 9.2 & 26\,483 & 0 & 27/30\\
MLA & 98.6 & 95.5 & 0 & 9.2 & 24\,667 & 0 & 27/30\\
MCA & 98.6 & 95.5 & 0 & 9.2 & 23\,607 & 0 & 27/30\\
LP & 94.2 & 73.8 & 0 & 17.5 & 24\,995 & 0 & 7/30\\
LT & 94.8 & 77.8 & 0 & 16.0 & 25\,793 & 0 & 7/30\\
DD & 93.6 & 71.0 & 0 & 29.0 & 28\,364 & 0 & 7/30\\
TD & 97.3 & 90.0 & 0 & 8.5 & 24\,818 & 0 & 27/30\\
MP & 96.3 & 83.0 & 0 & 15.0 & 26\,325 & 0 & 27/30\\
HP & 95.6 & 85.3 & 0 & 12.8 & 26\,857 & 0 & 27/30\\
MS & 95.3 & 82.8 & 0 & 14.5 & 27\,477 & 0 & 27/30\\
OF & 95.2 & 81.8 & 0 & 13.8 & 25\,793 & 0 & 7/30\\
TP & 96.7 & 84.9 & 0 & 13.0 & 25\,261 & 0 & 27/30\\
TT & 97.9 & 92.0 & 0 & 6.5 & 24\,464 & 0 & 27/30\\
TU & 92.0 & 65.8 & 0 & 20.4 & 29\,870 & 0 & 7/30\\
TG & 94.1 & 77.8 & 0 & 16.2 & 27\,743 & 0 & 27/30\\
TM & 95.0 & 82.8 & 0 & 14.1 & 26\,059 & 0 & 27/30\\
RG & 97.5 & 88.9 & 0 & 11.8 & 21\,849 & 0 & 27/30\\
RM & 97.6 & 89.9 & 0 & 11.5 & 21\,598 & 0 & 27/30\\
HG & 95.8 & 88.8 & 0 & 10.5 & 25\,120 & 0 & 27/30\\
HM & 95.8 & 88.8 & 0 & 10.7 & 25\,424 & 0 & 27/30\\
\bottomrule\end{tabular*}
\end{minipage}\par\medskip
\clearpage
\noindent\begin{minipage}{\textwidth}\singlespacing\captionsetup{font=footnotesize}
\setlength{\tabcolsep}{4pt}
\renewcommand{\arraystretch}{1.0}
\captionof{table}{\illustrativecaption Core-30 Combined. Energy and economic outcomes.}
\label{tab:full-scorecard-4-0}
\fontsize{8}{9}\selectfont
\begin{tabular*}{\textwidth}{@{\extracolsep{\fill}}lrrrrrrrr@{}}
\toprule
Code & Cost & CO$_2$ & Import & Export & Solar & Daily & Max. & Ramp\\
 & [EUR] & [kg] & [kWh] & [kWh] & [\%] & [kW] & [kW] & [kW]\\
\midrule
B0 & 25\,500 & 33\,900 & 189\,000 & 39\,200 & 30.0 & 60.0 & 100.0 & 186\,000\\
B1 & 23\,200 & 31\,200 & 174\,000 & 25\,800 & 53.9 & 57.0 & 97.0 & 178\,000\\
RB & 26\,504 & 30\,030 & 169\,167 & 25\,063 & 55.2 & 54.1 & 92.4 & 172\,660\\
RSL & 22\,968 & 27\,690 & 160\,467 & 21\,377 & 61.8 & 51.3 & 88.7 & 165\,540\\
RC & 22\,458 & 27\,008 & 157\,567 & 20\,640 & 63.1 & 49.0 & 84.1 & 156\,640\\
MLF & 22\,633 & 26\,817 & 163\,538 & 19\,309 & 65.5 & 51.0 & 86.7 & 165\,139\\
MCF & 22\,127 & 26\,218 & 159\,883 & 18\,878 & 66.3 & 49.9 & 84.8 & 161\,448\\
MLA & 20\,610 & 24\,420 & 148\,920 & 17\,583 & 68.6 & 46.5 & 79.0 & 150\,377\\
MCA & 19\,725 & 23\,372 & 142\,524 & 16\,829 & 69.9 & 44.5 & 75.6 & 143\,920\\
LP & 25\,399 & 28\,470 & 163\,367 & 23\,589 & 57.9 & 53.6 & 91.5 & 170\,880\\
LT & 24\,847 & 28\,080 & 161\,433 & 22\,851 & 59.2 & 52.4 & 89.6 & 167\,320\\
DD & 24\,019 & 27\,105 & 158\,050 & 21\,377 & 61.8 & 50.2 & 86.8 & 160\,200\\
TD & 22\,639 & 24\,570 & 149\,833 & 17\,691 & 68.4 & 46.7 & 79.4 & 151\,300\\
MP & 23\,743 & 26\,228 & 155\,633 & 19\,903 & 64.5 & 49.0 & 85.0 & 158\,420\\
HP & 23\,467 & 25\,838 & 154\,183 & 19\,166 & 65.8 & 48.5 & 83.1 & 156\,640\\
MS & 23\,743 & 26\,422 & 156\,600 & 20\,640 & 63.1 & 49.6 & 85.9 & 160\,200\\
OF & 24\,433 & 28\,080 & 161\,433 & 22\,851 & 59.2 & 52.4 & 89.6 & 167\,320\\
TP & 22\,713 & 27\,398 & 159\,017 & 22\,114 & 60.5 & 51.3 & 87.8 & 165\,540\\
TT & 21\,437 & 25\,058 & 151\,767 & 18\,429 & 67.1 & 47.9 & 81.3 & 154\,860\\
TU & 25\,123 & 27\,267 & 166\,278 & 19\,633 & 64.9 & 51.9 & 88.1 & 167\,906\\
TG & 24\,019 & 26\,069 & 158\,969 & 18\,770 & 66.5 & 49.6 & 84.3 & 160\,525\\
TM & 23\,467 & 25\,469 & 155\,315 & 18\,339 & 67.3 & 48.5 & 82.3 & 156\,836\\
RG & 22\,202 & 26\,767 & 155\,118 & 19\,903 & 64.5 & 48.5 & 85.7 & 156\,640\\
RM & 21\,947 & 26\,460 & 153\,335 & 19\,166 & 65.8 & 47.3 & 84.7 & 154\,860\\
HG & 22\,915 & 24\,869 & 151\,660 & 16\,954 & 69.7 & 44.5 & 80.4 & 142\,400\\
HM & 23\,191 & 25\,170 & 153\,487 & 16\,217 & 71.0 & 43.3 & 81.4 & 138\,840\\
\bottomrule\end{tabular*}
\end{minipage}\par\medskip
\noindent\begin{minipage}{\textwidth}\singlespacing\captionsetup{font=footnotesize}
\setlength{\tabcolsep}{4pt}
\renewcommand{\arraystretch}{1.0}
\captionof{table}{\illustrativecaption Core-30 Combined. Service, safety and resource use.}
\label{tab:full-scorecard-4-1}
\fontsize{8}{9}\selectfont
\begin{tabular*}{\textwidth}{@{\extracolsep{\fill}}lrrrrrrr@{}}
\toprule
Code & \gls{ev} min. & \gls{ev} band & Viol. & Corr. & Battery & \gls{v2g} & Saving\\
 & [\%] & [\%] & [kWh] & [\%] & [kWh] & [kWh] & [n/N]\\
\midrule
B0 & 93.0 & 16.0 & 0 & 19.0 & 0 & 0 & 0/30\\
B1 & 94.5 & 16.5 & 0 & 16.0 & 17\,500 & 0 & 0/30\\
RB & 92.2 & 62.4 & 0 & 20.0 & 18\,693 & 0 & 7/30\\
RSL & 96.3 & 85.7 & 0 & 17.0 & 20\,284 & 0 & 26/30\\
RC & 96.3 & 86.7 & 0 & 16.6 & 21\,318 & 0 & 26/30\\
MLF & 98.2 & 14.8 & 0 & 14.2 & 24\,310 & 0 & 26/30\\
MCF & 98.2 & 14.8 & 0 & 14.2 & 23\,767 & 0 & 26/30\\
MLA & 98.2 & 94.8 & 0 & 14.2 & 22\,137 & 0 & 26/30\\
MCA & 98.2 & 94.8 & 0 & 14.2 & 21\,186 & 0 & 26/30\\
LP & 91.2 & 68.4 & 0 & 22.5 & 22\,432 & 0 & 7/30\\
LT & 91.8 & 72.4 & 0 & 21.0 & 23\,148 & 0 & 7/30\\
DD & 91.6 & 71.4 & 0 & 27.0 & 25\,455 & 0 & 7/30\\
TD & 92.8 & 83.4 & 0 & 15.3 & 22\,273 & 0 & 26/30\\
MP & 92.5 & 78.4 & 0 & 18.5 & 23\,625 & 0 & 7/30\\
HP & 92.6 & 79.9 & 0 & 17.8 & 24\,102 & 0 & 7/30\\
MS & 92.3 & 77.4 & 0 & 19.5 & 24\,659 & 0 & 7/30\\
OF & 92.2 & 76.4 & 0 & 18.8 & 23\,148 & 0 & 7/30\\
TP & 95.5 & 82.7 & 0 & 18.0 & 22\,670 & 0 & 26/30\\
TT & 96.2 & 89.2 & 0 & 15.0 & 21\,955 & 0 & 26/30\\
TU & 89.0 & 60.4 & 0 & 25.4 & 26\,807 & 0 & 7/30\\
TG & 91.1 & 72.4 & 0 & 21.2 & 24\,898 & 0 & 7/30\\
TM & 92.0 & 77.4 & 0 & 19.1 & 23\,386 & 0 & 7/30\\
RG & 96.3 & 86.7 & 0 & 16.8 & 19\,608 & 0 & 26/30\\
RM & 96.4 & 87.7 & 0 & 16.5 & 19\,383 & 0 & 26/30\\
HG & 92.8 & 83.4 & 0 & 15.5 & 22\,544 & 0 & 26/30\\
HM & 92.8 & 83.4 & 0 & 15.7 & 22\,816 & 0 & 26/30\\
\bottomrule\end{tabular*}
\end{minipage}\par\medskip
\clearpage
\noindent\begin{minipage}{\textwidth}\singlespacing\captionsetup{font=footnotesize}
\setlength{\tabcolsep}{4pt}
\renewcommand{\arraystretch}{1.0}
\captionof{table}{\illustrativecaption Premium-100 Normal. Energy and economic outcomes.}
\label{tab:full-scorecard-5-0}
\fontsize{8}{9}\selectfont
\begin{tabular*}{\textwidth}{@{\extracolsep{\fill}}lrrrrrrrr@{}}
\toprule
Code & Cost & CO$_2$ & Import & Export & Solar & Daily & Max. & Ramp\\
 & [EUR] & [kg] & [kWh] & [kWh] & [\%] & [kW] & [kW] & [kW]\\
\midrule
B0 & 111\,000 & 146\,000 & 780\,000 & 290\,000 & 38.3 & 255.0 & 490.0 & 750\,000\\
B1 & 96\,000 & 128\,000 & 682\,000 & 210\,000 & 55.3 & 230.0 & 445.0 & 690\,000\\
RB & 92\,160 & 123\,200 & 663\,056 & 204\,000 & 56.6 & 218.5 & 423.8 & 669\,300\\
RSL & 86\,400 & 113\,600 & 628\,956 & 174\,000 & 63.0 & 207.0 & 406.9 & 641\,700\\
RC & 84\,480 & 110\,800 & 617\,589 & 168\,000 & 64.3 & 197.8 & 385.7 & 607\,200\\
MLF & 85\,920 & 110\,020 & 640\,993 & 157\,170 & 66.6 & 205.8 & 398.0 & 640\,144\\
MCF & 84\,000 & 107\,560 & 626\,667 & 153\,660 & 67.3 & 201.2 & 389.1 & 625\,838\\
MLA & 78\,240 & 100\,184 & 583\,697 & 143\,120 & 69.5 & 187.4 & 362.4 & 582\,924\\
MCA & 74\,880 & 95\,884 & 558\,630 & 136\,980 & 70.9 & 179.4 & 346.7 & 557\,892\\
LP & 88\,320 & 116\,800 & 640\,322 & 192\,000 & 59.1 & 216.2 & 419.6 & 662\,400\\
LT & 86\,400 & 115\,200 & 632\,744 & 186\,000 & 60.4 & 211.6 & 411.1 & 648\,600\\
DD & 83\,520 & 111\,200 & 619\,483 & 174\,000 & 63.0 & 202.4 & 398.4 & 621\,000\\
TD & 77\,400 & 100\,800 & 587\,278 & 144\,000 & 69.4 & 188.6 & 364.5 & 586\,500\\
MP & 82\,560 & 107\,600 & 610\,011 & 162\,000 & 65.5 & 197.8 & 389.9 & 614\,100\\
HP & 81\,600 & 106\,000 & 604\,328 & 156\,000 & 66.8 & 195.5 & 381.4 & 607\,200\\
MS & 82\,560 & 108\,400 & 613\,800 & 168\,000 & 64.3 & 200.1 & 394.1 & 621\,000\\
OF & 84\,960 & 115\,200 & 632\,744 & 186\,000 & 60.4 & 211.6 & 411.1 & 648\,600\\
TP & 85\,440 & 112\,400 & 623\,272 & 180\,000 & 61.7 & 207.0 & 402.6 & 641\,700\\
TT & 78\,300 & 102\,800 & 594\,856 & 150\,000 & 68.1 & 193.2 & 373.0 & 600\,300\\
TU & 87\,360 & 111\,864 & 651\,734 & 159\,800 & 66.0 & 209.3 & 404.3 & 650\,873\\
TG & 83\,520 & 106\,948 & 623\,087 & 152\,780 & 67.5 & 200.1 & 386.5 & 622\,261\\
TM & 81\,600 & 104\,488 & 608\,765 & 149\,270 & 68.2 & 195.5 & 377.6 & 607\,959\\
RG & 83\,520 & 109\,812 & 607\,992 & 162\,000 & 65.5 & 195.5 & 393.3 & 607\,200\\
RM & 82\,560 & 108\,552 & 601\,001 & 156\,000 & 66.8 & 190.9 & 388.6 & 600\,300\\
HG & 79\,680 & 102\,028 & 594\,439 & 138\,000 & 70.6 & 179.4 & 368.7 & 552\,000\\
HM & 80\,640 & 103\,260 & 601\,600 & 132\,000 & 71.9 & 174.8 & 373.4 & 538\,200\\
\bottomrule\end{tabular*}
\end{minipage}\par\medskip
\noindent\begin{minipage}{\textwidth}\singlespacing\captionsetup{font=footnotesize}
\setlength{\tabcolsep}{4pt}
\renewcommand{\arraystretch}{1.0}
\captionof{table}{\illustrativecaption Premium-100 Normal. Service, safety and resource use.}
\label{tab:full-scorecard-5-1}
\fontsize{8}{9}\selectfont
\begin{tabular*}{\textwidth}{@{\extracolsep{\fill}}lrrrrrrr@{}}
\toprule
Code & \gls{ev} min. & \gls{ev} band & Viol. & Corr. & Battery & \gls{v2g} & Saving\\
 & [\%] & [\%] & [kWh] & [\%] & [kWh] & [kWh] & [n/N]\\
\midrule
B0 & 99.0 & 18.0 & 0 & 6.0 & 0 & 0 & 0/100\\
B1 & 99.3 & 18.5 & 0 & 4.5 & 190\,000 & 0 & 0/100\\
RB & 98.7 & 74.1 & 0 & 9.0 & 202\,955 & 0 & 84/100\\
RSL & 99.3 & 91.1 & 0 & 6.0 & 220\,227 & 0 & 84/100\\
RC & 99.3 & 92.1 & 0 & 5.6 & 231\,455 & 0 & 84/100\\
MLF & 99.9 & 17.8 & 0 & 3.2 & 263\,936 & 0 & 84/100\\
MCF & 99.9 & 17.8 & 0 & 3.2 & 258\,037 & 0 & 84/100\\
MLA & 99.9 & 97.8 & 0 & 3.2 & 240\,341 & 0 & 84/100\\
MCA & 99.9 & 97.8 & 0 & 3.2 & 230\,021 & 0 & 84/100\\
LP & 97.7 & 80.1 & 0 & 11.5 & 243\,545 & 0 & 84/100\\
LT & 98.3 & 84.1 & 0 & 10.0 & 251\,318 & 0 & 84/100\\
DD & 98.1 & 83.1 & 0 & 16.0 & 276\,364 & 0 & 84/100\\
TD & 99.3 & 95.1 & 0 & 4.3 & 241\,818 & 0 & 84/100\\
MP & 99.0 & 90.1 & 0 & 7.5 & 256\,500 & 0 & 84/100\\
HP & 99.1 & 91.6 & 0 & 6.8 & 261\,682 & 0 & 84/100\\
MS & 98.8 & 89.1 & 0 & 8.5 & 267\,727 & 0 & 84/100\\
OF & 98.7 & 88.1 & 0 & 7.8 & 251\,318 & 0 & 84/100\\
TP & 98.5 & 88.1 & 0 & 7.0 & 246\,136 & 0 & 84/100\\
TT & 99.2 & 94.6 & 0 & 4.0 & 238\,364 & 0 & 84/100\\
TU & 95.5 & 72.1 & 0 & 14.4 & 291\,045 & 0 & 84/100\\
TG & 97.6 & 84.1 & 0 & 10.2 & 270\,318 & 0 & 84/100\\
TM & 98.5 & 89.1 & 0 & 8.1 & 253\,909 & 0 & 84/100\\
RG & 99.3 & 92.1 & 0 & 5.8 & 212\,886 & 0 & 84/100\\
RM & 99.4 & 93.1 & 0 & 5.5 & 210\,442 & 0 & 84/100\\
HG & 99.3 & 95.1 & 0 & 4.5 & 244\,763 & 0 & 84/100\\
HM & 99.3 & 95.1 & 0 & 4.7 & 247\,717 & 0 & 84/100\\
\bottomrule\end{tabular*}
\end{minipage}\par\medskip
\clearpage
\noindent\begin{minipage}{\textwidth}\singlespacing\captionsetup{font=footnotesize}
\setlength{\tabcolsep}{4pt}
\renewcommand{\arraystretch}{1.0}
\captionof{table}{\illustrativecaption Premium-100 All-in. Energy and economic outcomes.}
\label{tab:full-scorecard-6-0}
\fontsize{8}{9}\selectfont
\begin{tabular*}{\textwidth}{@{\extracolsep{\fill}}lrrrrrrrr@{}}
\toprule
Code & Cost & CO$_2$ & Import & Export & Solar & Daily & Max. & Ramp\\
 & [EUR] & [kg] & [kWh] & [kWh] & [\%] & [kW] & [kW] & [kW]\\
\midrule
B0 & 105\,000 & 140\,000 & 730\,000 & 264\,000 & 37.1 & 246.0 & 460.0 & 766\,000\\
B1 & 93\,000 & 125\,000 & 651\,000 & 198\,000 & 52.9 & 225.0 & 430.0 & 720\,000\\
RB & 109\,814 & 120\,312 & 632\,917 & 192\,343 & 54.2 & 213.8 & 409.5 & 698\,400\\
RSL & 94\,581 & 110\,938 & 600\,367 & 164\,057 & 60.9 & 202.5 & 393.1 & 669\,600\\
RC & 92\,479 & 108\,203 & 589\,517 & 158\,400 & 62.3 & 193.5 & 372.7 & 633\,600\\
MLF & 94\,056 & 107\,441 & 611\,857 & 148\,189 & 64.7 & 201.4 & 384.5 & 667\,976\\
MCF & 91\,954 & 105\,039 & 598\,182 & 144\,879 & 65.5 & 196.9 & 375.9 & 653\,048\\
MLA & 85\,648 & 97\,836 & 557\,166 & 134\,942 & 67.9 & 183.4 & 350.1 & 608\,268\\
MCA & 81\,970 & 93\,637 & 533\,238 & 129\,153 & 69.2 & 175.5 & 335.0 & 582\,148\\
LP & 105\,239 & 114\,062 & 611\,217 & 181\,029 & 56.9 & 211.5 & 405.4 & 691\,200\\
LT & 102\,951 & 112\,500 & 603\,983 & 175\,371 & 58.2 & 207.0 & 397.2 & 676\,800\\
DD & 99\,519 & 108\,594 & 591\,325 & 164\,057 & 60.9 & 198.0 & 385.0 & 648\,000\\
TD & 93\,800 & 98\,438 & 560\,583 & 135\,771 & 67.7 & 184.5 & 352.2 & 612\,000\\
MP & 98\,375 & 105\,078 & 582\,283 & 152\,743 & 63.6 & 193.5 & 376.8 & 640\,800\\
HP & 97\,232 & 103\,516 & 576\,858 & 147\,086 & 65.0 & 191.2 & 368.6 & 633\,600\\
MS & 98\,375 & 105\,859 & 585\,900 & 158\,400 & 62.3 & 195.8 & 380.9 & 648\,000\\
OF & 101\,235 & 112\,500 & 603\,983 & 175\,371 & 58.2 & 207.0 & 397.2 & 676\,800\\
TP & 93\,530 & 109\,766 & 594\,942 & 169\,714 & 59.6 & 202.5 & 389.0 & 669\,600\\
TT & 88\,276 & 100\,391 & 567\,817 & 141\,429 & 66.3 & 189.0 & 360.4 & 626\,400\\
TU & 104\,095 & 109\,242 & 622\,110 & 150\,669 & 64.1 & 204.8 & 390.7 & 679\,172\\
TG & 99\,519 & 104\,441 & 594\,764 & 144\,050 & 65.7 & 195.8 & 373.5 & 649\,316\\
TM & 97\,232 & 102\,039 & 581\,093 & 140\,740 & 66.5 & 191.2 & 364.9 & 634\,392\\
RG & 91\,428 & 107\,238 & 580\,356 & 152\,743 & 63.6 & 191.2 & 380.0 & 633\,600\\
RM & 90\,377 & 106\,008 & 573\,683 & 147\,086 & 65.0 & 186.8 & 375.5 & 626\,400\\
HG & 94\,944 & 99\,637 & 567\,419 & 130\,114 & 69.0 & 175.5 & 356.3 & 576\,000\\
HM & 96\,088 & 100\,840 & 574\,254 & 124\,457 & 70.4 & 171.0 & 360.8 & 561\,600\\
\bottomrule\end{tabular*}
\end{minipage}\par\medskip
\noindent\begin{minipage}{\textwidth}\singlespacing\captionsetup{font=footnotesize}
\setlength{\tabcolsep}{4pt}
\renewcommand{\arraystretch}{1.0}
\captionof{table}{\illustrativecaption Premium-100 All-in. Service, safety and resource use.}
\label{tab:full-scorecard-6-1}
\fontsize{8}{9}\selectfont
\begin{tabular*}{\textwidth}{@{\extracolsep{\fill}}lrrrrrrr@{}}
\toprule
Code & \gls{ev} min. & \gls{ev} band & Viol. & Corr. & Battery & \gls{v2g} & Saving\\
 & [\%] & [\%] & [kWh] & [\%] & [kWh] & [kWh] & [n/N]\\
\midrule
B0 & 91.0 & 16.0 & 0 & 23.0 & 0 & 0 & 0/100\\
B1 & 93.0 & 16.5 & 0 & 19.0 & 150\,000 & 0 & 0/100\\
RB & 90.2 & 58.8 & 0 & 23.0 & 160\,227 & 0 & 22/100\\
RSL & 94.8 & 83.0 & 0 & 20.0 & 173\,864 & 0 & 22/100\\
RC & 94.8 & 84.0 & 0 & 19.6 & 182\,727 & 0 & 82/100\\
MLF & 97.2 & 13.0 & 0 & 17.2 & 208\,370 & 0 & 22/100\\
MCF & 97.2 & 13.0 & 0 & 17.2 & 203\,714 & 0 & 82/100\\
MLA & 97.2 & 93.0 & 0 & 17.2 & 189\,743 & 0 & 82/100\\
MCA & 97.2 & 93.0 & 0 & 17.2 & 181\,595 & 0 & 82/100\\
LP & 89.2 & 64.8 & 0 & 25.5 & 192\,273 & 0 & 22/100\\
LT & 89.8 & 68.8 & 0 & 24.0 & 198\,409 & 0 & 22/100\\
DD & 89.6 & 67.8 & 0 & 30.0 & 218\,182 & 0 & 22/100\\
TD & 90.8 & 79.8 & 0 & 18.3 & 190\,909 & 0 & 22/100\\
MP & 90.5 & 74.8 & 0 & 21.5 & 202\,500 & 0 & 22/100\\
HP & 90.6 & 76.3 & 0 & 20.8 & 206\,591 & 0 & 22/100\\
MS & 90.3 & 73.8 & 0 & 22.5 & 211\,364 & 0 & 22/100\\
OF & 90.2 & 72.8 & 0 & 21.8 & 198\,409 & 0 & 22/100\\
TP & 94.0 & 80.0 & 0 & 21.0 & 194\,318 & 0 & 22/100\\
TT & 94.7 & 86.5 & 0 & 18.0 & 188\,182 & 0 & 82/100\\
TU & 87.0 & 56.8 & 0 & 28.4 & 229\,773 & 0 & 22/100\\
TG & 89.1 & 68.8 & 0 & 24.2 & 213\,409 & 0 & 22/100\\
TM & 90.0 & 73.8 & 0 & 22.1 & 200\,455 & 0 & 22/100\\
RG & 94.8 & 84.0 & 0 & 19.8 & 168\,068 & 0 & 82/100\\
RM & 94.9 & 85.0 & 0 & 19.5 & 166\,139 & 0 & 82/100\\
HG & 90.8 & 79.8 & 0 & 18.5 & 193\,234 & 0 & 22/100\\
HM & 90.8 & 79.8 & 0 & 18.7 & 195\,566 & 0 & 22/100\\
\bottomrule\end{tabular*}
\end{minipage}\par\medskip

